% (celles-ci ne doivent pas être jointes)
% (Dans la liste des auteurs, mettre votre nom en gras et souligner le nom des étudiants encadrés)
% Présentation des publications selon les spécificités disciplinaires. Les candidats sont invités à se reporter aux éventuelles préconisations formulées par leur section.

Comme demandé, je mets en gras mon nom et je souligne ceux des étudiants en master ou en thèse (mais pas en post-doc) que je co-encadrais lors de la publication. 
J'ai continué à collaborer avec d'anciens doctorants après leur thèse : je ne souligne alors plus leurs noms.
Les règles gouvernant l'ordre des auteurs varient selon le contexte et les disciplines (ordre alphabétique en mathématique ; premier, puis dernier, puis deuxième, puis avant-dernier, etc., en biologie ; hybride dans les travaux inter-disciplinaires). 
Les doctorants sont généralement les premiers auteurs des publications associées à leur thèse.

%---------------------------------------------------------------------
% \paragraph{Articles dans revues internationales à comité de lecture}
\subbib{../Biblio/Article}{Articles dans des revues internationales à comité de lecture}

\subbib{../Biblio/Preprint}{Pré-publications}

%---------------------------------------------------------------------
% \paragraph{Articles dans revues nationales à comité de lecture}
% \subbib{../Biblio/ArticleFr}{Articles dans des revues nationales à comité de lecture}

%---------------------------------------------------------------------
% \paragraph{Ouvrages individuels et direction d'ouvrages collectifs}
% \subbib{../Biblio/Livres-Rec}{Ouvrages destinés à la recherche}

% \subbib{../Biblio/Livres-Ens}{Ouvrages destinés à l'enseignement}

%---------------------------------------------------------------------
% \paragraph{Chapitres d'ouvrages}
% \subbib{../Biblio/Chapitre}{Chapitres d'ouvrage en anglais}

% \subbib{../Biblio/ChapitreFr}{Chapitres d'ouvrage en français}

%---------------------------------------------------------------------
\paragraph{Brevets, licences, logiciels.} 
J'ai été associé au développement de différents packages développé en langage R. 
Le seul publié depuis 2021 et disponible sur le site officiel du CRAN ({\sl The Comprehensive R Archive Network} : \url{cran.r-project.org/}) est le package \url{RGMM} : Robust Mixture Model associé à l'article \cite{GoR24}.
% \begin{itemize}
% %  \item \url{CGHseg} : Segmentation methods for array CGH analysis \cite{PLH11}
% %  \item \url{TAHMMAnnot} : HMM for ChIP and tiling-arrays using annotation \cite{BMB11}
% %  \item \url{HMMmix} : HMM with mixture as emission distribution \cite{VBM13}
% %  \item \url{Segmentor3IsBack} : A Fast Segmentation Algorithm \cite{CKL14}
% %  \item \url{EBS} : Exact Bayesian Segmentation \cite{ClR14}
%  \item \url{HiCseg} : Detection of domains in HiC data \cite{LDM14}
% %  \item \url{AR1seg} : Segmentation of an autoregressive Gaussian process of order 1 \cite{CLL14}
%  \item \url{SegCorr} : Detecting Correlated Genomic Regions \cite{DLM17}
%  \item \url{PLNmodels} : Poisson Lognormal Models \cite{CMR18a,CMR19,CMR21}
%  \item \url{RGMM} : Robust Mixture Model \cite{GoR24}
% \end{itemize}
% Ma contribution effective au développement de ces packages est variable et, pour tout dire, de plus en plus symbolique. 
% D'autres packages auquel j'ai contribué ne sont plus maintenus (\url{CGHseg}, \url{kerfdr}, \url{TAHMMAnnot}, \url{HMMmix}, \url{EBS}, \url{AR1seg}) (les unités de recherches ne sont pas toujours bien armées pour assurer la pérennité des outils qu'elles développent.). 
% J'ai été également associé au développement d'autres packages disponibles depuis github.

%---------------------------------------------------------------------
%\paragraph{Actes publiés de conférences internationales, congrès et colloques...}

%---------------------------------------------------------------------
%\paragraph{Autres}

